%% kapa: exam=1 round=21 year=2010 period=1 form=A subject=economics subject_ko=경제원론 seq=9 number=49
%% source: Q-Net 감정평가사 기출문제 (cst00309) artlSeq=5024181 file=extracted/1008/5024181/2010년 제21회 감정평가사 1차시험 기출문제/1교시.hwp
%% conversion: kordoc 4.24.51 -> tools/gen1.py (auto); review=proofread-sonnet; flags=none
\begin{question}{49}
A국의 생산함수는 $Y = AK^{\alpha}L^{\beta}$이다. 다음 자료를 바탕으로 성장회계에 의한 총요소생산성의 경제성장 기여율을 계산하면 얼마인가? (단, Y는 총소득, A는 총요소생산성, K는 자본스톡, L은 노동, $\alpha$는 자본소득분배율, $\beta$는 노동소득분배율이다.)
\begin{framedtext}
ㆍ연간 경제성장률 : 5\% \par
ㆍ연간 자본스톡증가율 : 7\% \par
ㆍ연간 노동증가율 : 1\% \par
ㆍ$\alpha = 0.5$ \par
ㆍ$\beta = 0.5$ \par
\end{framedtext}
\begin{choices}
  \item 10\%
  \item 15\%
  \item 20\%
  \item 25\%
  \item 30\%
\end{choices}
\end{question}
\answer{3}
